\subsection{极值点与重心}

命题 3. --- 设 K 是 Hausdorff 局部凸空间 E 的一个紧凸子集，x 是 K 的一个点。K 上每个以 x 为重心、总质量为 1 的正测度 \(\mu\)，都属于以 x 为重心、总质量为 1 的离散正测度所成之集的淡闭包。

设 U 是 \(\mu\) 对淡拓扑的一个邻域；我们可以假定 U 由 K 上满足下列条件的测度 \(\nu\) 组成：

\[
\left| \mu (f _ {i}) - \nu (f _ {i}) \right| \leqslant \delta\tag{1}
\]

其中 \(f_{i} \in \mathcal{C}(\mathrm{K}; \mathbf{C})\) ( \(1 \leqslant i \leqslant p\) ) 是有限个函数，\(\delta > 0\) 是一个数。对每个点 \(a \in K\)，存在 E 中 0 的一个闭凸邻域 \(V_{a}\)，使得

\[
\left| f _ {i} (\mathbf {y}) - f _ {i} (\mathbf {a}) \right| \leqslant \delta / 2\tag{2}
\]

对 \(1 \leqslant i \leqslant p\) 及每个 \(\mathbf{y} \in \mathrm{W}_{\mathbf{a}} = \mathrm{K} \cap (\mathbf{a} + \mathrm{V}_{\mathbf{a}})\) 成立。由于 \(\mathbf{K}\) 是紧的，存在有限个点 \(\mathbf{a}_j\) ( \(1 \leqslant j \leqslant r\) )（属于 \(\mathbf{K}\)），使得诸 \(\mathrm{W}_{\mathbf{a}_j}\) 构成 \(\mathbf{K}\) 的一个覆盖 ( \(1 \leqslant j \leqslant r\) )。取一个连续单位分解 \((g_j)_{1 \leqslant j \leqslant r}\)（在 \(\mathbf{K}\) 上，从属于覆盖 \((\mathrm{W}_{\mathbf{a}_j})\)），并令 \(\alpha_j = \mu(g_j)\)（对一切 \(j\)）；若 \(\alpha_j \neq 0\)，令 \(\mu_j = \alpha_j^{-1} g_j \cdot \mu\)，若 \(\alpha_j = 0\)，令 \(\mu_j = \varepsilon_{\mathbf{a}_j}\)。每个测度 \(\mu_j\) 都是正的、总质量为 1，且其支集含于紧凸集 \(\mathrm{W}_{\mathbf{a}_j}\)；此外，按定义，

\[
\mu = \sum_ {j = 1} ^ {r} \alpha_ {j} \mu_ {j}\tag{3}
\]

因为 \(g_{j} \cdot \mu = 0\)（若 \(\mu(g_{j}) = 0\)）；诸 \(\alpha_{j}\) 均 \(\geqslant 0\)，且

\[
\sum_ {j = 1} ^ {r} \alpha_ {j} = \sum_ {j = 1} ^ {r} \mu (g _ {j}) = \mu \left(\sum_ {j = 1} ^ {r} g _ {j}\right) = \mu (1) = 1.
\]

设 \(\mathbf{x}_j\) 是 \(\mu_j\) 的重心，它属于 \(\mathrm{W}_{\mathbf{a}_j}\)（No.~1, 命题 1），并考虑离散测度 \(\nu = \sum_{j=1}^{r} \alpha_j \varepsilon_{\mathbf{x}_j}\)；它是正的且总质量为 1，其重心是 \(\sum_{j=1}^{r} \alpha_j \mathbf{x}_j\)，由 (3)，后者也是 \(\mu\) 的重心，从而等于 \(\mathbf{x}\)。此外，由 (2)，有 \(|f_i(\mathbf{y}) - f_i(\mathbf{a}_j)| \leqslant \delta/2\)（对一切 \(\mathbf{y} \in \mathrm{W}_{\mathbf{a}_j}\) 及一切 \(i\)）；于是，由于 \(\operatorname{Supp}(\mu_j) \subset \mathrm{W}_{\mathbf{a}_j}\)，有 \(|\mu_j(f_i) - f_i(\mathbf{a}_j)| \leqslant \delta/2\)（对 \(1 \leqslant i \leqslant p\)）。另一方面，由于 \(\mathbf{x}_j \in \mathrm{W}_{\mathbf{a}_j}\)，还有

\[
\left| \varepsilon_ {\mathbf {x} _ {j}} (f _ {i}) - f _ {i} (\mathbf {a} _ {j}) \right| \leqslant \delta / 2
\]

对 \(1 \leqslant i \leqslant p\) 成立；从而 \(|\mu_j(f_i) - \varepsilon_{\mathbf{x}_j}(f_i)| \leqslant \delta\)（对一切 \(i\) 和 \(j\)）。由于诸 \(\alpha_j\) 均 \(\geqslant 0\) 且其和为 1，由 (3) 及 \(\nu\) 的定义推出 \(\nu\) 满足不等式 (1)。

证毕。

推论. --- 设 \(K'\) 是 K 的一个紧子集，且 K 是 \(K'\) 的闭凸包络。为使 \(x \in K'\) 是 K 的极值点，必须且只须 \(\varepsilon_{x}\) 是 \(K'\) 上以 x 为重心、总质量为 1 的唯一正测度。

设 \(\mathbf{x}\) 是 \(K\) 的一个极值点；为证明 \(\varepsilon_{\mathbf{x}}\) 是 \(K'\) 上以 \(\mathbf{x}\) 为重心、总质量为 1 的唯一正测度，由命题 3，只须看到离散测度 \(\nu\)（在 \(K'\) 上，为正的、总质量为 1，且以 \(\mathbf{x}\) 为重心）所成之集化为 \(\varepsilon_{\mathbf{x}}\)。但这样的测度 \(\nu\) 形如 \(\sum_{i=1}^{r} \lambda_i \varepsilon_{\mathbf{x}_i}\)，其中 \(\lambda_i > 0\)（对 \(1 \leqslant i \leqslant r\)）且 \(\sum_{i=1}^{r} \lambda_i = 1\)，而 \(\mathbf{x}\) 是 \(\nu\) 的重心这一假设可以写成 \(\mathbf{x} = \sum_{i=1}^{r} \lambda_i \mathbf{x}_i\)。由于 \(\mathbf{x}\) 是极值点，这就蕴涵 \(\mathbf{x}_i = \mathbf{x}\)（对一切 \(i\)），从而 \(\nu = \varepsilon_{\mathbf{x}}\)。

反之，假定 \(\varepsilon_{x}\) 是 \(K'\) 上以 x 为重心、总质量为 1 的唯一正测度，我们来证明 x 是极值点。若不然，就会存在 K 的两个不同点 \(x', x''\) 以及一个实数 \(\lambda\)，使得 \(0 < \lambda < 1\) 且 \(\mathbf{x} = \lambda\mathbf{x}' + (1 - \lambda)\mathbf{x}''\)。由命题 1，\(x'\)（相应地，\(x''\)）是某个正测度 \(\mu'\)（相应地，\(\mu''\)）的重心，该测度在 \(K'\) 上且总质量为 1。于是 x 是 \(\lambda\mu' + (1 - \lambda)\mu''\) 的重心。因此 \(\lambda\mu' + (1 - \lambda)\mu'' = \varepsilon_{x}\)。从而 \(\mu'\) 与 \(\mu''\) 都与 \(\varepsilon_{x}\) 成比例，于是 \(x' = x'' = x\)，这是荒谬的。

定理 1 (Choquet). --- 设 E 是 R 上的一个 Hausdorff 局部凸空间，K 是 E 的一个可度量化紧凸子集，M 是 K 的极值点所成之集。集 M 是 K 中一族可数开集的交，且 K 的每个点都是 K 上某个测度 \(\mu\) 的重心，其中 \(\mu(K - M) = 0\)。

为证明第一个断言，用 I 表示区间 {[}0,1{]}（\(\mathbf{R}\) 中的）；由于 K 是紧且可度量化的，\(\mathrm{K} \times \mathrm{K} \times \mathrm{I}\) 亦然；\(\mathrm{K} \times \mathrm{K} \times \mathrm{I}\) 中由三元组 \((\mathbf{x}, \mathbf{y}, \lambda)\)（满足 \(\mathbf{x} \neq \mathbf{y}\) 且 \(0 < \lambda < 1\)）组成的子集 U 在 \(\mathrm{K} \times \mathrm{K} \times \mathrm{I}\) 中是开集，因此存在一列闭集 \((\mathrm{F}_n)\)，它们在 \(\mathrm{K} \times \mathrm{K} \times \mathrm{I}\) 中，其并为 U（GT, IX, §2, No.~5, 命题 7）。映射 \(q: \mathrm{K} \times \mathrm{K} \times \mathrm{I} \to \mathrm{K}\)（由 \(q(\mathbf{x}, \mathbf{y}, \lambda) = \lambda \mathbf{x} + (1 - \lambda) \mathbf{y}\) 定义）是连续的，且由极值点的定义，\(\mathrm{K} - \mathrm{M} = q(\mathrm{U}) = \bigcup_{n} q(\mathrm{F}_n)\)；但 \(\mathrm{F}_n\) 是紧的，因为它在 \(\mathrm{K} \times \mathrm{K} \times \mathrm{I}\) 中是闭集，因此 \(q(\mathrm{F}_n)\) 是紧的，从而在 K 中是闭的；于是集 \(\mathrm{U}_n = \mathrm{K} - q(\mathrm{F}_n)\) 在 K 中是开集，且有 \(\mathrm{M} = \bigcap \mathrm{U}_n\)。

在证明的后续部分，我们将用 u 表示 K 上的一个连续凸数值函数，用 \(G \subset E \times R\) 表示 u 的图象，用 S 表示 G 在 \(E \times R\) 中的闭凸包络。

引理 1. --- 设 \(\overline{u}\) 是 E 上在 K 上 \(\geqslant u\) 的连续仿射线性函数的下包络。那么，S 是由点 \((\mathbf{a}, b) \in \mathrm{E} \times \mathbf{R}\)（满足 \(\mathbf{a} \in \mathrm{K}\) 且 \(u(\mathbf{a}) \leqslant b \leqslant \overline{u}(\mathbf{a})\)）组成之集。

由 Hahn-Banach 定理，为使 \((\mathbf{a}, b)\) 属于 S，必须且只须 \(h(\mathbf{a}, b) \geqslant 0\) 对 \(E \times R\) 上每个满足 \(h(\mathbf{x}, u(\mathbf{x})) \geqslant 0\)（对 \(x \in K\)）的连续仿射线性函数 h 成立。令 \(h(\mathbf{x}, t) = f(\mathbf{x}) - \lambda t\)，其中 f 是 E 上的一个连续仿射线性函数，我们看到关系 \((\mathbf{a}, b) \in S\) 等价于下列性质：关系

\[
f (\mathbf {x}) \geqslant \lambda u (\mathbf {x}) \quad \text { for   all } \mathbf {x} \in K\tag{4}
\]

蕴涵

\[
f (\mathbf {a}) \geqslant \lambda b.\tag{5}
\]

先设 \(\lambda = 0\)；由 Hahn-Banach 定理，对每个连续仿射线性函数 \(f\)（\(\mathbf{E}\) 上的），(4) 蕴涵 (5) 这一事实等价于关系 \(\mathbf{a} \in \mathbf{K}\)。其次，设 \(\lambda = -1\) 并把 \(f\) 换成 \(-f\)；那么关系 \(f(\mathbf{x}) \leqslant u(\mathbf{x})\)（在 \(\mathbf{K}\) 中）必须蕴涵 \(f(\mathbf{a}) \leqslant b\)。但由于 \(u\) 在 \(\mathbf{K}\) 上是凸的且连续的，\(u(\mathbf{a})\) 等于 \(f(\mathbf{a})\)（对函数 \(f\)——\(\mathbf{E}\) 上满足 \(f(\mathbf{x}) \leqslant u(\mathbf{x})\)（在 \(\mathbf{K}\) 上）的连续仿射线性函数——而言）的上确界（TVS, II, §5, No.~4, 命题 5）；由此得到关系 \(b \geqslant u(\mathbf{a})\)。最后，设 \(\lambda = 1\)；这时说 (4) 蕴涵 (5)，按定义即表示 \(b \leqslant \overline{u}(\mathbf{a})\)。这就证明了引理，因为关系 (4)（相应地，(5)）等价于两边同乘一个标量 \(>0\) 所得的关系。

引理 2. --- 若 u 在 K 上是严格凸的，则对 K 的每个非极值点 x，有 \(u(\mathbf{x}) < \overline{u}(\mathbf{x})\)。

因为这时存在 K 的两个不同点 y,z，使得 \(\mathbf{x} = (\mathbf{y} + \mathbf{z})/2\)，于是由 u 的严格凸性有 \(u(\mathbf{x}) < (u(\mathbf{y}) + u(\mathbf{z}))/2\)。若 f 是 E 上的一个仿射线性函数，则 \(f(\mathbf{x}) = (f(\mathbf{y}) + f(\mathbf{z}))/2\)；把这一关系应用于 E 上在 K 上 \(\geqslant u\) 的连续仿射线性函数 f，便得

\[
\overline {{{u}}} (\mathbf {x}) \geqslant \left(\overline {{{u}}} (\mathbf {y}) + \overline {{{u}}} (\mathbf {z})\right) / 2 \geqslant \left(u (\mathbf {y}) + u (\mathbf {z})\right) / 2 > u (\mathbf {x}).
\]

这两个引理既已建立，于是（引理 1）有 \((\mathbf{a},\overline{u} (\mathbf{a}))\in S\)（对一切 \(\mathbf{a}\in K\)）。由于 G 是 K 在连续映射 \(\mathbf{x}\mapsto (\mathbf{x},u(\mathbf{x}))\) 下的象，故是紧的，由 No.~1 的命题 1，存在一个正测度 \(\nu\)（在 G 上、总质量为 1、以 \((\mathbf{a},\overline{u} (\mathbf{a}))\) 为重心）。由于投影 \(\mathrm{pr}_1\) 在 G 上的限制是 G 到 K 上的一个同胚，测度 \(\nu\) 可借助这一同胚转移，从而给出 K 上的一个测度 \(\mu\)（正的且总质量为 1），使得

\[
\mathbf {a} = \int \mathbf {x} d \mu (\mathbf {x}) \quad \mathrm{and} \quad \overline {{{u}}} (\mathbf {a}) = \int u (\mathbf {x}) d \mu (\mathbf {x}).\tag{6}
\]

关系 (6) 中的第一个表明 a 是 \(\mu\) 的重心。函数 \(\overline{u}\) 在 K 上是上半连续且有界的，因而是 \(\mu\)-可积的（\(\S4\), No.~4, 命题 5 的推论 1）；此外，由于函数 \(-\overline{u}\) 按定义是凸的，由 No.~1 命题 2 的推论，有

\[
\overline {{{u}}} (\mathbf {a}) \geqslant \int \overline {{{u}}} (\mathbf {x}) d \mu (\mathbf {x}),\tag{7}
\]

于是，与公式 (6) 中的第二个相比较，便得

\[
\int u (\mathbf {x}) d \mu (\mathbf {x}) \geqslant \int \overline {{{u}}} (\mathbf {x}) d \mu (\mathbf {x}).\tag{8}
\]

但由于在 K 上 \(u(\mathbf{x}) \leqslant \overline{u}(\mathbf{x})\)，关系 (8) 蕴涵 \(u(\mathbf{x}) = \overline{u}(\mathbf{x})\) 对 \(\mu\) 而言几乎处处成立。考虑到引理 2 可见，只要证明了下面的引理，定理 1 即告完成：

引理 3. --- 设 E 是 R 上的一个 Hausdorff 局部凸空间，K 是 E 的一个可度量化紧凸子集。那么，K 上存在一个严格凸的数值函数。

因为 Banach 空间 \(\mathcal{C}(\mathrm{K};\mathbf{R})\) 是可分的（GT, X, §3, No.~3, 定理 1），故子空间 \(\mathcal{A}\)（由 E 上连续仿射线性函数在 K 上的限制组成，含于 \(\mathcal{C}(\mathrm{K};\mathbf{R})\)）也是可分的。于是设 \((h_n)\) 是 \(\mathcal{A}\) 中一个稠密序列，并设 \(\alpha_{n} > \sup_{\mathbf{x}\in \mathbb{K}}|h_n(\mathbf{x})|\)。那么每个函数 \(h_n^2 / n^2\alpha_n^2\) 在 K 中都是凸的（TVS, II, §2, No.~8, 例），且以 \(h_n^2 / n^2\alpha_n^2\) 为一般项的级数是正规收敛的，因此其和 \(u\) 在 K 中是连续且凸的。剩下要证 \(u\) 是严格凸的，为此只须证明：对 K 的任意两个不同点 \(\mathbf{x},\mathbf{x}'\)，存在一个整数 \(n\)，使得 \(h_n^2\) 在以 \(\mathbf{x},\mathbf{x}'\) 为端点的线段上的限制是严格凸的；而为此只须 \(h_n(\mathbf{x})\neq h_n(\mathbf{x}')\)（同前）。现在，存在一个函数 \(h\in \mathcal{A}\) 使得 \(h(\mathbf{x})\neq h(\mathbf{x}')\)（TVS, II, §4, No.~1, 命题 2 的推论 1），又由于序列 \((h_n)\) 在 \(\mathcal{A}\) 中稠密，故存在一个 \(n\) 使得 \(h_n(\mathbf{x})\neq h_n(\mathbf{x}')\)。

证毕。

推论. --- 设 E 是 R 上的一个 Hausdorff 局部凸空间，C 是 E 中以 0 为顶点的真凸锥，它对 E 的弱化一致结构所诱导的一致结构是完备且可度量化的。设 M 是 C 的极值母线之并。对每个 \(x \in C\)，存在 C 的一个紧凸子集 K 以及 K 上一个总质量为 1 的测度 \(\lambda \geqslant 0\)，使得 \(K - (M \cap K)\) 是 \(\lambda\)-可忽略的，且 \(\lambda\) 的重心等于 x。

因为 x 属于 C 的某个帽 K（TVS, II, §7, No.~2, 命题 5），且 M∩K 包含 K 的极值点所成之集（同前，命题 4 的推论 1）。于是只须应用定理 1。

